\begin{enunciado}{\ejercicio}
  Calcular el determinante de $A$ en cada uno de los siguientes casos:
  \begin{enumerate}[label=(\alph*)]
    \begin{multicols}{2}
      \item $A =
        \matriz{cccc}{
          1 & 0 & 0 & 2 \\
          0 & -2 & 3 & -1 \\
          -1 & 0 & 1 & 4 \\
          0 & 1 & -2 & 0
        }$

      \item $A =
        \matriz{ccc}{
          i & 0 & 2 + i \\
          -1 & 1-i & 0 \\
          2 & 0 & -1
        }$
    \end{multicols}
  \end{enumerate}
\end{enunciado}

\begin{enumerate}[label=(\alph*)]
  \item\label{ej-18:item-a} Acomodo un poco para ver si hago menos cuentas:
        $$
          \matriz{cccc}{
            1 & 0 & 0 & 2 \\
            0 & -2 & 3 & -1 \\
            -1 & 0 & 1 & 4 \\
            0 & 1 & -2 & 0
          }
          \triangulacion{
            F_3 + F_1 \to F_3\\
            F_4 + \frac{1}{2}\cdot F_2 \to F_4
          }
          \matriz{cccc}{
            1 & 0 & 0 & 2 \\
            0 & -2 & 3 & -1 \\
            0 & 0 & 1 & 6 \\
            0 & 0 & -\frac{1}{2} & -\frac{1}{2}
          }
          \triangulacion{
            F_4 + \frac{1}{2} \cdot F_3 \to F_4
          }
          \matriz{cccc}{
            1 & 0 & 0 & 2 \\
            0 & -2 & 3 & -1 \\
            0 & 0 & 1 & 6 \\
            0 & 0 & 0 & \frac{5}{2}
          }
        $$
        Ahora calculo el determinante, multiplicando los elementos de la diagonal porque es triangular:
        $$
          \textstyle
          \deter{cccc}{
            \blue{1} & 0         & 0        & 2        \\
            0        & \blue{-2} & 3        & -1       \\
            0        & 0         & \blue{1} & 6        \\
            0        & 0         & 0        & \blue{5}
          }
          = 1 \cdot (-2) \cdot 1 \cdot (-\frac{5}{2}) = \cajaResultado{ -5 }
        $$
        \parrafoDestacado[\atencion]{
          Este ejercicio inicialmente lo hice mal. Cuando se hace la operación de triangulación si se multiplica
          alguna \magenta{fila} por algún \orange{escalar}, esa fila no debe ser la que se modifica, por ejemplo en el primer paso:
          $$
            \ub{F_4}{\text{La que se}\\ \text{modifica}} + \orange{\frac{1}{2}}\cdot \ub{\magenta{F_2}}{\text{La que no se}\\ \text{modifica}} \to F_4
          $$
          porque sino se estaría \textit{escalando} la matriz así cambiando el determinante. Cuidado.
        }

  \item  Desarrollo por la segunda columna:
        $$
          \matriz{ccc}{
            i & 0 & 2 + i \\
            -1 & 1-i & 0 \\
            2 & 0 & -1
          }
          = (-1)^4 \cdot (1-i)
          \cdot
          \deter{cc}{
            i & 2+i \\
            2 & -1
          }
          = (1-i) \cdot (-4 - 3i) = \cajaResultado{-7 + i}
        $$
\end{enumerate}

\begin{aportes}
  \item \aporte{\dirRepo}{naD GarRaz \github}
\end{aportes}
